\documentclass[11pt,a4paper]{article}
\usepackage{Qsim}



\newcommand{\myTIPTitle}{\includegraphics[width=8cm]{logo}}
\newcommand{\myTIPSubtitle}{Simulador de Arquitecturas Q}

\title{\myTIPTitle \\\myTIPSubtitle\\ Manual de Instalación}
\author{Organización de Computadoras \and TPI}
\date{}


\newcommand{\lnx}{GNU/Linux}

\newcommand{\zip}{qsim.zip}

\pagestyle{fancy}
\fancyhead{} % limpio la cabecera\
\fancyhead[L]{\includegraphics[width=2cm]{logo}}
\fancyhead[R]{}

%\fancyhead[L]{\nouppercase{\leftmark \sectiontitle}}
\fancyfoot[C]{\thepage}
\fancyfoot[R]{}


\begin{document}

\maketitle
\tableofcontents

\section{Pre-requisitos}

\subsection{Para ejecutar en  \lnx}

Para ejecutar el simulador en el sistema operativo \lnx\ es necesario tener instalada la versión 7 del Java Runtime Environment. Para instalarla:
\begin{enumerate}
\item Abrir una terminal con la combinación de teclas \code{Ctrl+Alt+t}
\item Ejecutar \underline{en la terminal} el siguiente comando (que se puede pegar en la terminal usando el mouse):
\begin{verbatim}
sudo apt-get install openjdk-7-jre
\end{verbatim}
\item Ingresar la contraseña del usuario
\end{enumerate}

\subsection{Para ejecutar en Windows}

En construcción...

\section{Arranque del Simulador QSim}

El archivo comprimido \code{\zip} incluye dos archivos ejecutables, uno para cada posible sistema operativo, a saber:

\begin{itemize}
\item El script de arranque \code{qsim.sh} para \lnx\
\item El script de arranque \code{qsim.bat} para Windows
\end{itemize}



%--------------------------------------------------------------------------------------------

\subsection{Ejecución de \qsim\ en \lnx}

El script de arranque \code{qsim.sh} está preparado para una arquitectura de 32 bits. 

\begin{enumerate}
\item Abrir una terminal con la combinación de teclas \code{Ctrl+Alt+t}
\item Ubicarse en la carpeta que contiene los archivos descomprimidos. Si se lo descomprimió en el escritorio, la ruta es \code{Escritorio/ejecutable\_QSim} (dentro de la carpeta home del usuario). \\Para eso \textbf{ingresar el siguiente comando en la terminal}:
\begin{verbatim}
cd Escritorio/ejecutable_QSim
\end{verbatim}
\item Cambiar los permisos de ejecución sobre el script ingresando el siguiente comando en la terminal
\begin{verbatim}
sudo chmod u+x qsim.sh
\end{verbatim}
\item \textbf{Si lo pide}, ingresar la contraseña del usuario
\item Ejecutar el script, ingresando el siguiente comando en la terminal
\begin{verbatim}
sh qsim.sh
\end{verbatim}
\end{enumerate}

\subsection{Ejecución de \qsim\ en Windows}

 
\begin{enumerate}
\item Abrir la carpeta que contiene los archivos descomprimidos.
\item Hacer doble clic sobre el script de arranque \code{qsim.bat}
\end{enumerate}


\graf{qsim-pantallaPpal}{Pantalla inicial de \qsim}

\end{document}